\documentclass[10pt,a4paper]{amsart}
\usepackage[margin=19mm]{geometry}
\usepackage{amsmath,amssymb,mathtools}
\usepackage[scr=rsfs,cal=euler]{mathalfa}
\usepackage{xfrac}
\usepackage[unicode,hidelinks,bookmarks=false]{hyperref}
\newcommand{\ie}{i.e.\ }
\newcommand{\cf}{cf.\ }
\newcommand{\resp}{respectively\ }
\newcommand{\Subsection}{\subsection}
\providecommand{\nobreakdash}{\nobreak\hbox{-}}
\setlength{\emergencystretch}{2em}
\multlinegap0pt
\usepackage{bm}
\usepackage{bbm}
\usepackage{dsfont}
\usepackage{xspace}
\usepackage{cancel}
\usepackage{mathtools}
\usepackage{amsthm}
\makeatletter
\@ifundefined{theorem}{\newtheorem{theorem}{Theorem}}{}
\@ifundefined{prop}{\newtheorem{prop}[theorem]{Proposition}}{}
\makeatother
\makeatletter
\newcommand{\C}{\mathbb{C}}
\newcommand{\HS}{\operatorname{HS}}
\newcommand{\N}{\mathbb{N}}
\expandafter\ifx\csname customcoro\endcsname\relax
\newenvironment{customcoro}[1]
  {\renewcommand\theinnercustomthm{#1}\innercustomcoro}
  {\endinnercustomthm}
\fi
\expandafter\ifx\csname customthm\endcsname\relax
\newenvironment{customthm}[1]
  {\renewcommand\theinnercustomthm{#1}\innercustomthm}
  {\endinnercustomthm}
\fi
\makeatother
\title[Directional Poincar\'e inequality on compact Lie groups]{Directional Poincar\'e inequality on compact Lie groups}
\author{Paulo L. Dattori da Silva, Andr\'e Pedroso Kowacs}
\date{Source: arXiv:2510.05409v1 (CC BY 4.0)}
\begin{document}
\maketitle

\noindent\emph{Source and licence.} Directional Poincar\'e inequality on compact Lie groups. Version arXiv:2510.05409v1 (CC BY 4.0), distributed under the Creative Commons Attribution 4.0 International licence, as recorded in the arXiv licence metadata for this exact version and shown on the abstract page https://arxiv.org/abs/2510.05409v1. Full rights notice: licenses/arxiv-2510.05409-CC-BY-4.0.txt.\par\smallskip

\noindent\textit{Editorial scope.} Counted unit: the Prop whose source label is \texttt{prop\_globally\_solvable} in the article (source0.tex lines 960--966, source label \texttt{prop\_globally\_solvable}), delivered together with the article's own complete original proof of it (source0.tex lines 967--1055, one \texttt{proof} environment of 89 lines). No mathematics is written, repaired or re-derived: statement and proof are the article's own text line for line. Required context retained uncounted: prop\_kernel\_annihi\_range (lines 903--906). Every definition, display and label of those lines is retained unchanged. Omitted from this package and not redistributed: 1--902, 907--959, 1056--1254. Nothing inside a retained excerpt is omitted or altered: every retained body is delivered exactly as the article's lines are written, so no omission receipt of any kind accompanies this package. Citations: the retained excerpts carry no citation command at all, so no bibliography entry of the article is transcribed and no reference is redistributed. Reference adaptation: none was needed and none was made; every reference inside the delivered unit resolves inside it, so no unresolved reference or citation is produced. Printed slips kept literal: no printed word, symbol, hypothesis or number of the counted statement or of its proof was altered. Layout and class: the article's own class is \texttt{amsart} with packages including inputenc, amsmath, amssymb, wrapfig, url, mathtools, graphicx, stmaryrd, amsthm, xcolor, hyperref, enumitem, comment, relsize; the wrapper is the lane's pinned \texttt{amsart} editorial wrapper at 10pt on A4, which restates the theorem-like environments the retained text needs and adds the article's own macro definitions that the retained text needs (\texttt{\textbackslash C}, \texttt{\textbackslash HS}, \texttt{\textbackslash N}), with extra wrapper packages (bm, bbm, dsfont, xspace, cancel, mathtools). The delivered numbering is the unit's own. Assets: no retained span contains a figure environment, a graphics-inclusion command, a PDF-page inclusion, a table or a TikZ picture; no image file is copied into this package, the package declares no asset dependency and no image file is redistributed with it. Licence: version arXiv:2510.05409v1 of this article is distributed under the Creative Commons Attribution 4.0 International licence, as recorded in the arXiv licence metadata for this exact version and shown on the abstract page https://arxiv.org/abs/2510.05409v1. A full rights notice is delivered at licenses/arxiv-2510.05409-CC-BY-4.0.txt. Characters: every retained excerpt is pure ASCII, so the delivered engine, XeLaTeX with its default Latin Modern text font, reports no missing character.
\begin{prop}\label{prop_kernel_annihi_range}
    Let $G$ be a compact Lie group and $P$ a left-invariant continuous linear operator on $C^\infty(G)$. Then $f\in C^\infty(G)$ is in $(\ker {}^tP)^{0}$ if and only if, for every $[\xi]\in\widehat{G}$, 
    every column of $\widehat{f}(\xi)$ is in  $\operatorname{ran}\sigma_P(\xi)$.
\end{prop}

\begin{prop}\label{prop_globally_solvable}
    Let $G$ be a compact Lie group and $P:C^\infty(G)\to C^\infty(G)$ be a left-invariant continuous linear operator on $G$. Then $P$ is globally solvable (has closed range)  if and only if there exist $C,k>0$ such that
    \begin{equation}\label{ineq_diof_gs}
        \lambda_{\min}^{>0}[\sigma_P(\xi)]\geq C\langle \xi\rangle^{-k},
    \end{equation}
    for every $[\xi]\in\widehat{G}$ such that $\sigma_P(\xi)\neq 0$.
\end{prop}

\begin{proof}
Suppose that the inequality \eqref{ineq_diof_gs} is satisfied, and let $f\in C^\infty(G)\cap (\ker {}^tP)^0$.  We will construct $u\in C^\infty(G)$ such that $Pu=f$, as follows.

Let $[\xi]\in\widehat{G}$. If $\sigma_P(\xi)=0$, we must have $\widehat{f}(\xi)=0$. Then, let $\widehat{u}(\xi)=0$. 

Now assume $\sigma_P(\xi)\neq 0$. Due to the decomposition 
\begin{equation*}
    \ker \sigma_P(\xi)\oplus( \ker \sigma_P(\xi))^\perp=\C^{d_\xi},
\end{equation*}
by conjugating $\xi$ with a unitary change of basis (which amounts to choosing a different representative of the class $[\xi])$, we may assume that $\sigma_P(\xi)$ is given in block form
\begin{equation*}
    \sigma_P(\xi)=\begin{pmatrix}
        \tilde \sigma_P(\xi)&0\\
        0&0
    \end{pmatrix},
\end{equation*}
where the matrix $\tilde \sigma_P(\xi)$ corresponds to the restriction
\begin{equation*}
    \tilde \sigma_P(\xi)= \sigma_P(\xi)|_{(\ker \sigma_P(\xi))^\perp}:(\ker \sigma_P(\xi))^\perp\to \operatorname{ran}\sigma_P(\xi),
\end{equation*}
is invertible and satisfies
\begin{equation*}
     \lambda_{\min}[\tilde\sigma_P(\xi)]= \lambda_{\min}^{>0}[\sigma_P(\xi)]\geq C\langle \xi \rangle^{-k}.
\end{equation*}
This also implies that 
\begin{equation*}
    \| \tilde \sigma_P(\xi)^{-1}\|_{\operatorname{op}}\leq C^{-1}\langle \xi \rangle^{k},
\end{equation*}
and since the columns of $\widehat{f}(\xi)$ are in the range of $\sigma_P(\xi)$, by Proposition \ref{prop_kernel_annihi_range}, we must have that $\widehat{f}(\xi)$ is given in block form by
\begin{equation*}
    \widehat{f}(\xi)= \begin{pmatrix}
        {\widehat{f_1}}(\xi)\\0
    \end{pmatrix},
\end{equation*}
where $\widehat{f_1}(\xi)\in\C^{d^0_\xi\times d_\xi}$, $d^0_\xi = \dim \operatorname{ran}\sigma_P(\xi)=\dim (\ker \sigma_P(\xi))^{\perp}$.

Therefore, defining
\begin{equation*}
    \widehat{u}(\xi)=\begin{pmatrix}
        \tilde \sigma_P(\xi)^{-1}\widehat{f_1}(\xi)\\
        0
    \end{pmatrix}\in\C^{d_\xi\times d_\xi},
\end{equation*}
we have that $\sigma_P(\xi)\widehat{u}(\xi)=\widehat{f}(\xi)$ and, moreover,
\begin{align*}
    \|\widehat{u}(\xi)\|_{\HS}&=\| \tilde \sigma_P(\xi)^{-1}\widehat{f_1}(\xi)\|_{\HS}\\
    &\leq \| \tilde \sigma_P(\xi)^{-1}\|_{\operatorname{op}}\|\widehat{f}(\xi)\|_{\HS}\\
    &\leq C^{-1}\langle \xi \rangle^{k}\|\widehat{f}(\xi)\|_{\HS}.
\end{align*}
Due to the rapid decay of the Fourier coefficients of $f$, it follows that $u\in C^\infty(G)$. Moreover, the equality $\sigma_P(\xi)\widehat{u}(\xi)=\widehat{f}(\xi)$ for every $[\xi]\in\widehat{G}$ guarantees that $Pu=f$. This proves that $P$ is globally solvable.

Conversely, assume that inequality \eqref{ineq_diof_gs} does not hold. Then there exist sequences of distinct elements $[\xi_n]\in\widehat{G}$ and $v_{n}\in (\ker \sigma_P(\xi_n))^\perp\subset \C^{d_{\xi_n}}$, $\|v_n\|_2=1$, such that
\begin{equation*}
  0< \| \sigma_P(\xi_n)v_n\|_2<\langle \xi_n \rangle ^{-n},
\end{equation*}
for every $n\in\N$.

Consider $f\in C^\infty(G)$ given by the Fourier coefficients:
\begin{equation*}
    \widehat{f}(\xi)=\begin{cases}
        0&\text{ if }\xi\neq \xi_n,\ \forall n\in\N,\\
        \begin{pmatrix}
            \sigma_P(\xi_n)v_n&0&\cdots&0 
        \end{pmatrix}_{d_{\xi_n}\times d_{\xi_n}}&\text{ if }\xi= \xi_n,\ n\in\N.
    \end{cases}
\end{equation*}
Then, $f$ is well-defined, since 
\begin{equation*}
    \|\widehat{f}(\xi_n)\|_{\HS}=\|\sigma_P(\xi_n)v_n\|_2<\langle \xi_n \rangle ^{-n},
\end{equation*}
for every $n\in\N$ and $ \|\widehat{f}(\xi)\|_{\HS}=0$, for every other $[\xi]\in\widehat{G}$. Moreover, note that the columns of $\widehat{f}(\xi)$ are in $ \operatorname{ran}\sigma_P(\xi)$, for every $[\xi]\in\widehat{G}$, since 
\begin{equation*}
    \sigma_P(\xi_n)\begin{pmatrix}
            v_n&0&\cdots&0 
        \end{pmatrix}_{d_{\xi_n}\times d_{\xi_n}}=\widehat{f}(\xi_n),
\end{equation*}
for every $n\in\N$. Therefore by Proposition \ref{prop_kernel_annihi_range}, $f\in (\ker {}^tP)^0$. On the other hand, if $Pu=f$, then  the Fourier coefficients of $u$ must satisfy
\begin{equation*}
    \widehat{u}(\xi_n)=
        \begin{pmatrix}
            v_n&0&\cdots&0 
        \end{pmatrix}_{d_{\xi_n}\times d_{\xi_n}}+w_{\xi_n},
\end{equation*}
where the columns of $w_{\xi_n}\in \C^{d_\xi\times d_\xi}$ are in $\ker \sigma_P(\xi_n)$, and hence $\sigma_P(\xi_n)w_{\xi_n}=0$. Therefore, by orthogonality, we conclude that 
\begin{equation*}
    \|\widehat{u}(\xi_n)\|_{\HS}\geq \|v_n\|_2=1,
\end{equation*}
for every $n\in\N$. Therefore,  $u\in\mathcal{D}'(G)\backslash C^\infty(G)$ and, consequently, $P$ is not globally solvable.
\end{proof}

\end{document}
